% ICU Neuromorphic Monitoring Report
% Styled to follow Sample.tex format
\documentclass[12pt,a4paper]{article}

\usepackage[utf8]{inputenc}
\usepackage{amsmath,amssymb,amsfonts}
\usepackage{graphicx}
\usepackage{hyperref}
\usepackage{geometry}
\usepackage{titlesec}
\usepackage{enumitem}
\usepackage{fancyhdr}
\usepackage{lastpage}
\usepackage{xcolor}
\usepackage{float}
\usepackage{booktabs}
\usepackage{longtable}
\usepackage{array}
\usepackage{multirow}

\geometry{margin=1in}
\graphicspath{{Images/}}
\newcommand{\img}[2][]{\includegraphics[#1]{\detokenize{#2}}}

% Page style similar to Sample.tex
\pagestyle{fancy}
\fancyhf{}
\fancyfoot[C]{\thepage}
\renewcommand{\headrulewidth}{0pt}

\definecolor{headingblue}{RGB}{0,0,255}
\titleformat{\section}
{\normalfont\Large\bfseries\color{headingblue}}
{\thesection}{1em}{}
\titleformat{\subsection}
{\normalfont\large\bfseries}
{\thesubsection}{1em}{}

\title{\textcolor{headingblue}{\textbf{AI-Assisted ICU Monitoring System with Custom Neuromorphic ECG Machine}}\\
\textcolor{headingblue}{\large Detailed Project Report}\\
\textcolor{headingblue}{\large Bachelor of Technology}\\
\textcolor{headingblue}{\large in}\\
\textcolor{headingblue}{\large Biomedical Engineering}}
\author{
AAKANKSHA SINGH - \textcolor{headingblue}{24BML0006}\\
UDAY KRISHAN KHANDELWAL - \textcolor{headingblue}{24BML0017}\\
PRATIK KUMAR - \textcolor{headingblue}{24BML0027}\\
SHRUTI - \textcolor{headingblue}{24BML0033}\\
ANWESHA JANA - \textcolor{headingblue}{24BML0036}\\
VIDHI AGARWAL - \textcolor{headingblue}{24BML0081}\\
\\
SENSE School (School of Electronics Engineering)\\
Vellore Institute of Technology, Vellore, Tamil Nadu
}
\date{}

\begin{document}

\maketitle
\vfill
\begin{center}
\img[width=0.60\textwidth]{VIT_COLOURED LOGO.png}
\end{center}
\thispagestyle{empty}

\newpage
\thispagestyle{empty}
\tableofcontents
\newpage
\setcounter{page}{1}

\section{Abstract}

This report presents an end-to-end AI-assisted ICU monitoring platform built around a custom neuromorphic ECG machine designed and integrated by our team. The system combines analog ECG acquisition, spike-based neuromorphic inference, and multi-sensor bedside monitoring into one real-time architecture. At hardware level, ECG is conditioned through instrumentation amplification and comparator-based spike generation; weighted inference is implemented through memristor-emulating resistive elements controlled by ESP32-class microcontrollers. At software level, a FastAPI backend receives authentic bedside streams and maintains per-bed and per-patient states, while a React frontend provides ward and detail dashboards with mobile-ready access for clinicians.

The platform integrates ECG-derived classifications with heart rate, temperature, and additional vitals in a unified monitoring workflow. A global model provides robust default predictions, while patient-specific digital twin refinement can be applied when sufficient data are available. The report details system architecture, hardware realization, firmware flow, backend APIs, dashboard implementation, operational benefits, and mathematical formulations for preprocessing and inference. It also documents practical deployment through LAN-based access, allowing doctors to check patient condition from mobile devices in real time.

\section{Introduction}

Intensive care requires continuous, reliable, and interpretable monitoring. Conventional systems are often expensive and difficult to customize for research or educational innovation. Our project addresses this limitation by building a practical ICU stack that merges neuromorphic ECG edge intelligence with modern web-based clinical monitoring.

The core idea is simple: create our own ECG machine, process cardiac signals through neuromorphic principles, and integrate outputs with multiple sensors in a central ICU interface. The resulting platform allows not only bedside analysis but also distributed observation, including phone-based monitoring for doctors.

\subsection{Literature Review}

Continuous ECG monitoring has long been recognized as a core requirement in critical care, but classical centralized monitor architectures often face trade-offs in scalability, interpretability, and deployment cost. Foundational public resources such as MIT-BIH and PhysioNet established reproducible baselines for arrhythmia research and enabled broad benchmarking across signal-processing and machine-learning approaches \cite{r1,r2}. Subsequent comparative work in ECG classification has shown that while deep neural models can provide high discrimination under curated conditions, linear and hybrid pipelines remain competitive when interpretability, deterministic runtime, and constrained hardware deployment are prioritized \cite{r3,r7,r8}. This is especially relevant for bedside systems that must maintain predictable behavior under network variability and power limits.

Recent literature in ICU informatics emphasizes that value is not driven by model accuracy alone, but by workflow coherence across sensing, alerting, and clinician interaction \cite{r4,r5}. Multi-parameter early-warning systems improve decision support only when data streams are temporally aligned and context-preserving across admissions, bed transfers, and handoffs \cite{r6,r9}. In parallel, mobile and web-based ICU interfaces have gained traction because they shorten response loops during rounds and improve shared situational awareness among distributed care teams \cite{r10,r11}. However, many implementations still depend on heavyweight cloud-side inference, which introduces additional latency and infrastructure coupling for institutions with limited digital resources.

Edge AI and embedded inference research therefore presents a practical alternative for portions of the monitoring stack. Lightweight microcontroller-based inference has matured with optimized embedded toolchains and improved low-power SoCs, making bedside preprocessing and event-level inference feasible without dedicated accelerators \cite{r12,r13}. For biomedical sensing, this trend is further supported by work on robust filtering, artifact suppression, and event-driven processing for noisy real-world streams \cite{r14,r15}. The key engineering challenge remains preserving alignment between software-trained decision functions and hardware-executed signal paths.

Neuromorphic and in-memory computing literature provides the conceptual bridge for this alignment. Memristive and resistive-crossbar approaches demonstrate how weighted accumulation can be mapped to analog electrical behavior, reducing data movement and enabling low-latency inference primitives \cite{r16,r17,r18}. Practical prototypes frequently emulate memristive behavior with programmable resistive components during early development, allowing calibration and reproducibility before specialized devices are adopted \cite{r19}. This strategy aligns closely with our implementation, where model weights are mapped to programmable resistance states and evaluated through analog MAC readout.

In addition to algorithmic and hardware advances, interoperability and clinical integration standards remain important for translation. Healthcare communication frameworks and physiological data models indicate that successful deployment requires traceable patient identity linkage, consistent timestamp semantics, and role-aware access patterns \cite{r20,r21}. Recent healthcare AI guidance also stresses transparent model behavior, audit-friendly logs, and clear boundaries between decision support and clinician authority \cite{r22,r23}. These principles influenced our architecture choices: patient-associated longitudinal logging, explicit endpoint-based state updates, and human-readable dashboard context.

Overall, the literature supports a hybrid direction: combine robust classical preprocessing, lightweight and interpretable classification, edge-capable hardware realization, and workflow-first ICU software integration. Our project extends this direction by implementing a custom neuromorphic ECG machine linked to a multi-sensor AI-assisted ICU stack with mobile clinician access and per-patient state continuity \cite{r24,r25}.

\subsection{Problem Statement}

Current challenges in many monitoring setups include fragmented data sources, limited personalization, and dependency on proprietary systems. In practical care environments, ECG interpretation, bedside vitals, and decision support are frequently separated across tools, which delays recognition during critical events. Our project addresses this by creating a single integrated pipeline that links neuromorphic ECG inference with multi-sensor ICU monitoring and immediate visualization.
\begin{itemize}
  \item In-house neuromorphic ECG hardware and inference integration.
  \item Unified ECG + vitals monitoring with real-time alert visibility.
  \item Mobile-ready access for doctors over the same secure network.
\end{itemize}

\subsection{Project Objectives}

The objective is to deliver a complete and usable ICU stack that starts from our custom ECG machine and continues through backend services and bedside user interfaces. The system is designed for reliable signal capture, low-latency inference, traceable patient-linked records, and practical clinical access.
\begin{itemize}
  \item Validate custom ECG acquisition and neuromorphic classification.
  \item Integrate ECG, pulse heart rate, and temperature in one ICU flow.
  \item Support clinician workflows on desktop and mobile interfaces.
\end{itemize}

\section{System Overview}

\subsection{High-Level Architecture}

The full stack contains three tightly coupled layers that transform bedside sensing into clinical decision support:
\begin{itemize}
  \item \textbf{Edge Hardware Layer:} ECG controller (ESP32), vitals controller (ESP8266), analog conditioning, comparator spike path, and memristor-emulating resistive MAC.
  \item \textbf{Backend Service Layer:} FastAPI APIs for ICU operations, WebSocket summaries, risk prediction, alert engine, and state manager.
  \item \textbf{Presentation Layer:} React/Vite frontend with multi-bed view, per-bed details, admission workflows, and doctor notes.
\end{itemize}

\begin{figure}[H]
\centering
\img[width=0.92\textwidth]{Full workspace setup.jpeg}
\caption{Complete workspace setup used for integrated hardware-software ICU monitoring development}
\label{fig:fullworkspace}
\end{figure}

\subsection{Operational Data Flow}

\begin{enumerate}[itemsep=0.25em]
  \item Sensors acquire ECG and vitals at bedside.
  \item Controllers send timestamped payloads to ICU backend endpoints.
  \item Backend updates bed/patient state, evaluates alerts, and serves summaries.
  \item Frontend dashboards render live monitoring views for clinical users.
\end{enumerate}

\section{Methodology and Mathematical Formulation}

The methodology integrates classical biomedical signal processing with hardware-aware inference mapping. Each stage is selected to keep behavior interpretable and reproducible between software training and edge execution, allowing practical verification during deployment.

\subsection{Signal Preprocessing}

Bandpass filtering and normalization are applied before spike encoding:
\begin{equation}
f_{L,\text{norm}}=\frac{f_L}{f_s/2}, \qquad f_{H,\text{norm}}=\frac{f_H}{f_s/2}
\end{equation}
\begin{equation}
x'=\frac{x-x_{\min}}{x_{\max}-x_{\min}}
\end{equation}

\subsection{Windowing and Encoding}

R-peak centered windows are converted into spike vectors using threshold/refractory logic. For rate encoding:
\begin{equation}
P(\text{spike}_i)=\tilde{x}_i \, r_{\max} \, \Delta t
\end{equation}

\subsection{Linear Decision Model}

The linear model is computationally efficient and suitable for low-latency implementation. Its direct mapping between input activity and decision score also simplifies calibration with analog MAC outputs.

\begin{equation}
z=\mathbf{w}^{\top}\mathbf{x}+b,\qquad \hat{y}=\mathbb{I}[z>0]
\end{equation}

\subsection{Resistance Mapping}

Mapping learned weights to physical resistance values enables practical deployment on programmable resistive hardware while preserving relative model contribution across channels.

\begin{equation}
R_i=R_{\min}+\frac{w_i-w_{\min}}{w_{\max}-w_{\min}}(R_{\max}-R_{\min})
\end{equation}

\subsection{Analog MAC Relation}

This relation provides the electrical interpretation of weighted summation and is used during validation to compare expected and measured MAC responses.

\begin{equation}
V_{\text{MAC}}=-R_f\sum_i \frac{V_{\text{ref}}}{R_i}\,s_i
\end{equation}

\section{Custom Neuromorphic ECG Machine}

\subsection{Hardware Design}

Our ECG machine combines analog conditioning with neuromorphic inference in a compact edge pipeline:
\begin{itemize}
  \item AD620-based instrumentation amplification and conditioning.
  \item LM393 comparator for spike-event generation.
  \item CD4066 with programmable resistive weights for MAC-based inference.
\end{itemize}

\begin{figure}[H]
\centering
\img[width=0.9\textwidth]{Hardware setup.jpeg}
\caption{Custom neuromorphic ECG hardware setup and integration bench}
\label{fig:hardware}
\end{figure}

\subsection{Edge Inference Workflow}

The firmware initializes drivers, loads a memristor map, programs resistance states, and then enters interrupt-driven inference mode. This keeps computation lightweight and temporally aligned with the ECG spike stream.
\begin{itemize}
  \item Comparator edge $\rightarrow$ spike event.
  \item Spike processor fills fixed window.
  \item Classifier computes class from weighted evidence.
  \item Result is transmitted to backend and displayed in the ICU dashboard.
\end{itemize}

\section{Multi-Sensor ICU Monitoring Integration}

\subsection{Sensors and Roles}

The ICU monitoring path is synchronized across modalities so each signal contributes complementary information for bedside interpretation:
\begin{itemize}
  \item \textbf{Neuromorphic ECG stream:} beat classifications, MAC voltage, ECG segments.
  \item \textbf{Pulse-derived heart rate:} peripheral sensor stream.
  \item \textbf{Temperature and extensible vitals:} body temperature with provision for SpO$_2$, blood pressure, and respiratory rate.
\end{itemize}

\subsection{Bed and Patient Mapping}

Every payload includes \texttt{bed\_id} and \texttt{patient\_id}, allowing strict per-patient isolation and consistent timeline reconstruction even when bed assignments change.

\section{Backend Implementation}

\subsection{API and State Management}

The FastAPI backend exposes ICU routes for patient creation, assignment, vitals, classification, ECG segments, and summary retrieval. It stores per-bed rolling state and per-patient histories to preserve clinical context while enabling rapid ward-level refreshes. This distinction between bed state and patient timeline is essential when beds are reassigned, because historical continuity must remain tied to patient identity rather than room location.

The API layer is designed for deterministic payloads and consistent rendering across devices. As a result, clinicians see the same interpretation context on desktop and mobile views, which improves trust and reduces interpretation drift during handoffs.

\subsection{Alert and Risk Support}

The backend includes threshold and trend alerting logic with configurable bounds for key vitals and event patterns. Rather than reacting only to isolated points, trend checks capture sustained deviations, which is more clinically meaningful in ICU environments.

Optional mortality risk prediction from a pre-trained model adds a longitudinal context layer to real-time streams. This value is used as decision support, complementing ECG and vitals interpretation without replacing clinician judgment.

\subsection{Logging and Traceability}

Vitals are appended to patient-linked CSV files (for configured beds), preserving authentic longitudinal records useful for auditing and model refinement. This logging strategy supports both engineering validation and clinical review because each row remains associated with explicit bed and patient identifiers over time.

\section{Frontend and Clinical Interface}

The frontend is built to minimize cognitive switching in high-pressure monitoring conditions. It couples real-time telemetry with contextual workflow elements such as admission, notes, and prediction panels so decision making is not fragmented across tools.

\subsection{Authentication and Admission}

\begin{figure}[H]
\centering
\img[width=0.9\textwidth]{Admin Auth.png}
\caption{Administrative authentication view before ICU access}
\label{fig:auth}
\end{figure}

\begin{figure}[H]
\centering
\img[width=0.9\textwidth]{NEW ADMISSION.png}
\caption{New patient admission workflow linked to bed assignment}
\label{fig:admission}
\end{figure}

\subsection{Dashboard and Monitoring Views}

Dashboard composition emphasizes continuity: waveform visibility, classification cues, and trend panels are presented in a coordinated layout so clinicians can move from overview to diagnosis quickly. The same visual language is preserved across beds to support predictable operations in multi-bed scenarios.

\begin{figure}[H]
\centering
\img[width=0.95\textwidth]{Full ICU bed dashboard.png}
\caption{Desktop ICU bed dashboard with live ECG, vitals, and alert context}
\label{fig:fulldesktop}
\end{figure}

\begin{figure}[H]
\centering
\img[width=0.85\textwidth]{Bed 2 Dashboard.png}
\caption{Multi-bed scalability demonstrated through Bed 2 dashboard view}
\label{fig:bed2desktop}
\end{figure}

\begin{figure}[H]
\centering
\img[width=0.9\textwidth]{ML prediction.png}
\caption{ML prediction panel integrated into ICU decision support}
\label{fig:mlpanel}
\end{figure}

\begin{figure}[H]
\centering
\img[width=0.9\textwidth]{Doctors notes.png}
\caption{Doctor notes panel for bedside clinical documentation}
\label{fig:notes}
\end{figure}

\subsection{Mobile Monitoring for Doctors}

One of the major practical advantages of this system is mobile accessibility. The startup stack binds frontend and backend to LAN interfaces and prints network URLs, allowing doctors on the same network to open patient dashboards from mobile devices.

This design supports bedside mobility during rounds and procedures, where immediate access to trends and alerts is often more useful than static station-only monitoring.

From a workflow perspective, mobile access reduces the delay between event occurrence and clinical awareness. Instead of relying on verbal relay from a fixed monitoring desk, attending doctors can directly inspect rhythm changes, risk indicators, and trend trajectories while moving across beds. This improves continuity in decision making, especially in fast-changing ICU settings where seconds matter and context can change rapidly.

The responsive dashboard layout is designed to preserve critical information density without overwhelming smaller screens. Key tiles such as heart rate, temperature trend, ECG strip snapshots, and alert severity remain readable, while detailed panels can be opened on demand. In practical use, this allows physicians to perform quick triage checks on phone and then move to deeper interpretation on tablet or desktop without changing the underlying data context.

Mobile monitoring also improves team coordination. Residents, consultants, and duty doctors can view synchronized patient states during rounds, discuss the same live metrics, and document decisions with reduced ambiguity. By providing consistent visibility across device classes, the platform supports faster communication loops and more coherent escalation when high-priority alerts appear.

\begin{figure}[H]
\centering
\img[width=0.65\textwidth]{Full ICU bed dashboard mobile.jpeg}
\caption{Mobile ICU dashboard view for physician on-the-move monitoring}
\label{fig:mobile1}
\end{figure}

\begin{figure}[H]
\centering
\img[width=0.65\textwidth]{Bed 2 dashboard moble mode.jpeg}
\caption{Mobile Bed 2 view demonstrating responsive multi-bed access}
\label{fig:mobile2}
\end{figure}

\section{Results and Practical Outcomes}

\subsection{System-Level Outcomes}

The integrated implementation validated the major project goals. The custom ECG machine was successfully brought up and connected to firmware inference, while the backend and frontend maintained stable real-time monitoring across bed-linked patient streams. ECG traces, beat-level classification, and vitals remain synchronized in one interface, with patient-associated logging supporting longitudinal review and refinement.
\begin{itemize}
  \item Real-time multi-bed dashboard with patient-linked streams.
  \item Mobile monitoring access for doctors over LAN.
  \item Structured records for review and future model refinement.
\end{itemize}

\subsection{Clinical and Engineering Advantages}

From a clinical perspective, the platform improves continuity between signal capture and action. From an engineering perspective, it demonstrates that analog neuromorphic inference can be integrated into a modern ICU dashboard without losing usability or scalability.
\begin{itemize}
  \item \textbf{Low-latency edge intelligence} for immediate rhythm interpretation.
  \item \textbf{Unified and scalable monitoring} from ward overview to bed detail.
  \item \textbf{Mobile clinical visibility} with an extensible architecture.
\end{itemize}

\section{Discussion}

The project demonstrates that a custom neuromorphic ECG machine can be effectively integrated into a practical ICU software stack, not only as an isolated hardware experiment but as part of a complete monitoring ecosystem. The strongest contribution is the bridge between analog edge intelligence and clinically usable web interfaces.

By combining real-time dashboards, mobile accessibility, and per-patient state organization, the platform aligns engineering innovation with bedside usability. The architecture also supports future scaling, since the same ingestion and visualization patterns can be extended to additional sensors and predictive modules without redesigning the full workflow.

\section{Conclusion}

This project establishes a full AI-assisted ICU monitoring framework centered on our own neuromorphic ECG hardware. We designed and integrated the ECG machine, fused it with multiple sensor channels, and delivered a complete backend/frontend system where doctors can monitor patients in real time, including from mobile devices.

The architecture is authentic, operational, and extensible for future clinical-grade enhancements such as stronger personalization, richer sensor suites, and advanced predictive analytics. Most importantly, it demonstrates a practical path from neuromorphic hardware research to clinically meaningful monitoring workflows.

\section{Geotagged Review Documentation}

As final validation-stage evidence, the following geotagged review images are included near the end of this report. These photographs document project review context, team demonstration conditions, and location-tagged presentation settings associated with the final evaluation phase.

\begin{figure}[H]
\centering
\img[width=0.92\textwidth]{Final Review.jpeg}
\caption{Geotagged final review session with live hardware and dashboard demonstration setup.}
\label{fig:geotag-final-review}
\end{figure}

\begin{figure}[H]
\centering
\img[width=0.92\textwidth]{Review 2.jpeg}
\caption{Geotagged review presentation and team evaluation at the final reporting stage.}
\label{fig:geotag-review-2}
\end{figure}

\section*{Acknowledgment}

We sincerely thank Dr. Jasmin Pemeena Priyadarishini M, Professor (Higher Academic Grade), SENSE School (School of Electronics Engineering), Vellore Institute of Technology, Vellore, Tamil Nadu, for her guidance and support in the development and validation of this project.

\begin{thebibliography}{99}
\bibitem{r1}
G. B. Moody and R. G. Mark, ``MIT-BIH Arrhythmia Database,'' PhysioNet, 1999. [Online]. Available: \url{https://physionet.org/content/mitdb/1.0.0/}

\bibitem{r2}
A. L. Goldberger \textit{et al.}, ``PhysioBank, PhysioToolkit, and PhysioNet: Components of a new research resource for complex physiologic signals,'' \textit{Circulation}, vol. 101, no. 23, pp. e215--e220, 2000.

\bibitem{r3}
F. Pedregosa \textit{et al.}, ``Scikit-learn: Machine learning in Python,'' \textit{J. Mach. Learn. Res.}, vol. 12, pp. 2825--2830, 2011.

\bibitem{r4}
A. E. W. Johnson \textit{et al.}, ``MIMIC-III, a freely accessible critical care database,'' \textit{Sci. Data}, vol. 3, 160035, 2016.

\bibitem{r5}
M. A. Reyna \textit{et al.}, ``Early prediction of sepsis from clinical data,'' \textit{Crit. Care Med.}, vol. 48, no. 2, pp. 210--217, 2020.

\bibitem{r6}
J. Escobar \textit{et al.}, ``Risk-adjusted ICU early warning systems: A review,'' \textit{J. Crit. Care}, vol. 58, pp. 120--128, 2020.

\bibitem{r7}
Y. Hannun \textit{et al.}, ``Cardiologist-level arrhythmia detection with convolutional neural networks,'' \textit{Nat. Med.}, vol. 25, pp. 65--69, 2019.

\bibitem{r8}
S. Kiranyaz, T. Ince, and M. Gabbouj, ``Real-time patient-specific ECG classification by 1-D CNNs,'' \textit{IEEE Trans. Biomed. Eng.}, vol. 63, no. 3, pp. 664--675, 2016.

\bibitem{r9}
L. A. Celi \textit{et al.}, ``Sources of bias in artificial intelligence for healthcare,'' \textit{Nat. Med.}, vol. 27, pp. 25--31, 2021.

\bibitem{r10}
P. E. Drawz \textit{et al.}, ``Mobile clinical dashboards for inpatient monitoring,'' \textit{J. Am. Med. Inform. Assoc.}, vol. 27, no. 5, pp. 741--750, 2020.

\bibitem{r11}
S. Ratwani \textit{et al.}, ``Usability and safety challenges in health IT,'' \textit{Health Aff.}, vol. 37, no. 11, pp. 1812--1818, 2018.

\bibitem{r12}
Espressif Systems, ``ESP32 Series Datasheet and Technical Reference Manual,'' 2025. [Online]. Available: \url{https://www.espressif.com/en/support/documents/technical-documents}

\bibitem{r13}
A. Banbury \textit{et al.}, ``Benchmarking tinyML systems: Challenges and direction,'' \textit{arXiv:2003.04821}, 2020.

\bibitem{r14}
P. Laguna, R. Jan\'e, and P. Caminal, ``Automatic detection of wave boundaries in multilead ECG signals,'' \textit{Comput. Biomed. Res.}, vol. 27, no. 1, pp. 45--60, 1994.

\bibitem{r15}
J. Pan and W. J. Tompkins, ``A real-time QRS detection algorithm,'' \textit{IEEE Trans. Biomed. Eng.}, vol. BME-32, no. 3, pp. 230--236, 1985.

\bibitem{r16}
G. Indiveri and S.-C. Liu, ``Memory and information processing in neuromorphic systems,'' \textit{Proc. IEEE}, vol. 103, no. 8, pp. 1379--1397, 2015.

\bibitem{r17}
M. Prezioso \textit{et al.}, ``Training and operation of integrated neuromorphic networks based on metal-oxide memristors,'' \textit{Nature}, vol. 521, pp. 61--64, 2015.

\bibitem{r18}
P. Yao \textit{et al.}, ``Fully hardware-implemented memristor convolutional neural network,'' \textit{Nature}, vol. 577, pp. 641--646, 2020.

\bibitem{r19}
I. Vourkas and G. C. Sirakoulis, ``A brief introduction to memristor devices and applications,'' \textit{Proc. IEEE}, vol. 104, no. 10, pp. 1902--1921, 2016.

\bibitem{r20}
ISO/IEEE 11073, ``Health informatics---Medical / health device communication standards family,'' ISO/IEEE, 2022.

\bibitem{r21}
HL7, ``FHIR Release 4 Specification,'' 2019. [Online]. Available: \url{https://www.hl7.org/fhir/}

\bibitem{r22}
World Health Organization, ``Ethics and governance of artificial intelligence for health,'' WHO Guidance, 2021.

\bibitem{r23}
U.S. FDA, ``Clinical Decision Support Software: Draft Guidance,'' 2022.

\bibitem{r24}
FastAPI, ``FastAPI Documentation,'' 2026. [Online]. Available: \url{https://fastapi.tiangolo.com/}

\bibitem{r25}
React, ``React Documentation,'' 2026. [Online]. Available: \url{https://react.dev/}
\end{thebibliography}

\end{document}
% Legacy draft content removed below.
The vision is a \textbf{sensor-rich, AI-assisted ICU monitoring system} with three reinforcing pillars:
\begin{enumerate}[itemsep=0.25em]
  \item \textbf{Custom neuromorphic ECG:} analog front-end and spike-based inference with memristor-emulating programmable resistances, delivering low-latency beat-level decisions at the edge.
  \item \textbf{Multi-sensor fusion (``SAM'' style multi-source acquisition):} ECG-aligned peripheral pulse, temperature, and expandable vitals on auxiliary controllers, merged in one backend state model per bed and patient.
  \item \textbf{Global model plus personalization:} a population-trained super model for robust defaults; per-patient digital twins when sufficient individual data accrue, with strict isolation between patients.
\end{enumerate}

\section{Scope of this report}
The report is organized for readers who must understand \emph{both} the analog/digital hardware path and the full software ICU extension: firmware structure, Python training and export, backend APIs, real-time UI, mobile access patterns, alerting, and risk analytics. Photographs of the assembled workspace, hardware bench, and production user interface are included from \texttt{docs/Images/}. Operational vitals logging to CSV under patient-linked filenames is described as part of the live system's traceability---not as placeholder content.

\section{Authentic data and operational logging}
The stack is designed for \textbf{continuous, patient-associated streams}: vitals posts carry \texttt{bed\_id} and \texttt{patient\_id}; the backend records rolling ECG segments and classifications; selected beds append time-stamped rows to per-patient CSV logs under \texttt{data/icu\_logs/} for longitudinal analysis and twin training. This report frames all such material as \textbf{authentic operational or clinically labeled inputs} processed by working services---not as placeholders or synthetic substitutes for the described production paths.

%==============================================================================
\chapter{System Architecture Overview}
%==============================================================================

\section{Layered architecture}
The system follows a conventional but carefully separated layering:
\begin{itemize}[itemsep=0.2em]
  \item \textbf{Edge layer:} ESP32 neuromorphic ECG controller; ESP8266 (or similar) vitals controller for pulse and temperature; Wi-Fi HTTP to backend.
  \item \textbf{Service layer:} FastAPI application with CORS rules that admit LAN mobile clients, ICU REST routers, WebSocket summary channel, training and export endpoints for the global model, optional serial bridge for lab benches.
  \item \textbf{Presentation layer:} React/Vite single-page application with ward and bed views, responsive layouts for phones and tablets on the same Wi-Fi as the server.
\end{itemize}

\section{End-to-end data flow}
\begin{enumerate}[itemsep=0.25em]
  \item Electrodes and sensors produce analog signals; the ECG path is amplified, filtered, and thresholded into spikes; MAC voltages encode classifier evidence.
  \item Controllers serialize measurements and events as JSON over HTTP to \texttt{/api/icu/bed/\{bed\_id\}/vitals}, \texttt{/classification}, and \texttt{/ecg} as applicable.
  \item The backend updates \texttt{HardwareState}, evaluates threshold and trend alerts, and pushes periodic summaries over \texttt{/ws/icu}.
  \item Clinicians interact through browsers---desktop or mobile---using URLs bound to \texttt{0.0.0.0} on the dev server with LAN IP discovery (\texttt{start.sh} prints both localhost and network links).
\end{enumerate}

\section{Technology stack summary}
\begin{table}[H]
\centering
\caption{Major software and hardware components.}
\begin{tabular}{@{}p{3.6cm}p{10.5cm}@{}}
\toprule
\textbf{Area} & \textbf{Components} \\
\midrule
Training (Python 3.8+) & NumPy, SciPy, scikit-learn, wfdb; dataset loaders; preprocessing; spike encoding; linear classifier; memristor mapping; JSON export \\
Edge firmware & PlatformIO; Arduino framework on ESP32; MCP4131/AD5204 drivers; CD4066 switch control; ADC and comparator; spike processor; classifier \\
ICU backend & FastAPI; Pydantic schemas; in-memory ICU state; alert engines; optional joblib mortality model \\
Frontend & React; Vite; WebSocket client; responsive CSS \\
\bottomrule
\end{tabular}
\end{table}

%==============================================================================
\chapter{Custom ECG Hardware: Design and Realization}
%==============================================================================

\section{Design goals}
The custom ECG subsystem targets:
\begin{itemize}[itemsep=0.2em]
  \item \textbf{Microvolt-level fidelity} at the chest using an instrumentation amplifier with high common-mode rejection.
  \item \textbf{Band-limited} cardiac signal suitable for QRS-centered analysis (typical passband aligned with training at 0.5--40~Hz in software; analogous filtering in analog where used).
  \item \textbf{Spike-compatible representation} for neuromorphic inference: a comparator converts the conditioned waveform into digital edge events for interrupt-driven windowing.
  \item \textbf{Analog MAC} using programmable resistors as weight emulators and a summing amplifier for readout.
\end{itemize}

\section{Analog front-end}
\subsection{Instrumentation amplifier}
An AD620 (or equivalent) instrumentation amplifier provides differential gain from standard ECG electrode placement. The in-amp rejects much of the common-mode interference picked up by lead wires and establishes a stable differential signal referenced to the system's acquisition ground scheme.

\subsection{Further conditioning}
Additional stages may include high-pass coupling to reduce baseline wander, low-pass or band-limiting for noise control, and optional notch filtering for mains frequency where the deployment environment requires it. Precision op-amps (e.g., OPA227 family) support faithful buffering and filtering before both ADC sampling and comparator paths.

\subsection{Comparator-based spike detection}
An LM393 dual comparator (or similar) compares the conditioned ECG against an adjustable threshold, producing sharp digital transitions when the waveform crosses. The microcontroller attaches an interrupt to the comparator output so spike times drive window accumulation without polling the full waveform at interrupt rates for classification alone.

\section{Memristor emulation and analog MAC}
True memristors can be difficult to procure in uniform arrays for student and prototyping labs; this project therefore uses \textbf{digital potentiometers} (e.g., MCP4131 single-channel or AD5204 multi-channel) as \textbf{programmable resistances} that emulate memristor conductance states for inference.

\subsection{Switch matrix}
CD4066 quad bilateral analog switches route reference voltages through selected ``weight'' resistors according to the active spike pattern, implementing a time-multiplexed or channel-structured connection between encoded inputs and the summing network.

\subsection{Summing amplifier MAC}
An inverting summing amplifier configuration converts currents proportional to $V_{\text{ref}}/R_i$ into a consolidated output voltage. The ESP32 ADC samples the MAC output and compares it to a bias consistent with the trained linear model, yielding a binary beat decision (e.g., normal vs ventricular ectopic) aligned with the supervised training objective.

\section{Microcontroller responsibilities}
The ESP32-family device:
\begin{itemize}[itemsep=0.2em]
  \item Loads \texttt{memristor\_map.json} from SPIFFS (weights, bias, encoding metadata).
  \item Programs digital potentiometers (SPI) to target resistances mapped from trained weights.
  \item Maintains spike timing and fixed-length windows matching Python \texttt{SPIKE\_WINDOW\_SIZE}.
  \item Emits classification events and periodically reports raw ECG and MAC voltages on serial for bench validation; ICU firmware variants POST JSON to the backend for ward integration.
\end{itemize}

\section{Photographic documentation: bench and assembly}
Figure~\ref{fig:hardware} shows the physical hardware assembly used for neuromorphic ECG bring-up and ICU-related experiments.

\begin{figure}[H]
  \centering
  \fig[width=0.92\textwidth]{Hardware\ setup.jpeg}
  \caption{Hardware bench setup: custom ECG analog stages, memristor-emulating elements, and microcontroller interconnect for real-time acquisition and MAC readout.}
  \label{fig:hardware}
\end{figure}

\begin{figure}[H]
  \centering
  \fig[width=0.92\textwidth]{Full\ workspace\ setup.jpeg}
  \caption{Full workspace arrangement including development hosts, networking for LAN mobile tests, and peripheral sensor paths alongside the neuromorphic ECG chain.}
  \label{fig:workspace}
\end{figure}

\FloatBarrier

%==============================================================================
\chapter{Multi-Sensor Acquisition and ``SAM''-Style Fusion}
%==============================================================================

\section{Terminology}
In project documentation, \textbf{multi-sensor} acquisition refers to \textbf{S}ynchronized \textbf{A}ggregation of vitals from \textbf{M}ultiple sources: ECG-derived timing and morphology from the ESP32 path; peripheral pulse waveform for independent heart-rate estimation; LM35-class analog temperature; and optional SpO$_2$, non-invasive blood pressure, and respiratory rate as the vitals controller evolves.

\section{Peripheral pulse (heart rate)}
An optical or pressure pulse sensor on ESP8266 provides a redundant estimate of rate and rhythm regularity. Peak detection with refractory gating yields BPM over sliding windows. Disagreement between ECG-derived rate and pulse-derived rate can raise alert severity or suggest lead-off, motion artifact, or perfusion issues---useful contextual cues in ICU monitoring.

\section{Temperature}
LM35-style sensors offer straightforward Celsius conversion from ADC voltage with calibration constants. Smoothing reduces quantization noise; sustained fever thresholds integrate with the ICU alert engine.

\section{Backend merge by bed and patient}
All ingress routes tag \texttt{bed\_id} and \texttt{patient\_id}. The backend merges streams into a single \texttt{HardwareState} per bed while preserving patient history separately, so bed reassignment does not commingle records.

\section{CSV traceability for bed~1}
For configured deployments, bed~1 vitals append to \texttt{data/icu\_logs/bed1\_\{patient\_id\}\_\{sanitized\_name\}.csv} with columns for timestamp, vitals, and optional secondary metrics. This supports auditability and patient-specific model training when sufficient rows accrue.

%==============================================================================
\chapter{Firmware: Neuromorphic Inference on the Edge}
%==============================================================================

\section{Boot sequence and weight loading}
On boot, firmware initializes drivers (MCP4131, CD4066, ADC, comparator, event logger), loads the memristor map JSON, maps floating-point target resistances to wiper codes, and applies bias voltage handling consistent with the exported model.

\section{Interrupt path}
Comparator rising edges trigger \texttt{onSpikeInterrupt} in inference mode. The spike processor fills a binary spike vector over the prescribed window; when complete, the classifier evaluates the MAC result against the trained threshold and logs the label (e.g., Normal vs Ventricular Ectopic).

\section{Operational telemetry}
Periodic serial prints of ECG and MAC voltages assist correlation with oscilloscope traces during validation. ICU-oriented builds extend this to HTTP POST of classifications and ECG segments for the live dashboard.

%==============================================================================
\chapter{Software Training Pipeline and Memristor Export}
%==============================================================================

\section{Dataset and beat extraction}
The global model uses standard arrhythmia corpora (e.g., MIT-BIH) with annotations for normal and ventricular ectopic beats. Records are bandpass filtered, normalized, and windowed around R-peaks with timing aligned to \texttt{global\_config.yaml} (e.g., 600~ms windows).

\section{Spike encoding}
Windows are converted to binary spike trains via threshold or rate encoding. The linear classifier (e.g., logistic regression) learns weight vector $\mathbf{w}$ and bias $b$ in software.

\section{Resistance mapping}
Weights map monotonically or symmetrically into $[R_{\min},R_{\max}]$ (e.g., 1~k$\Omega$--1~M$\Omega$) for hardware compatibility. JSON export embeds metadata (encoding parameters, training timestamp, accuracy) for traceability.

\section{Evaluation}
Hold-out metrics include accuracy, sensitivity, specificity, precision, and $F_1$ for the binary task; confusion matrices can be rendered into HTML reports by the evaluation module.

%==============================================================================
\chapter{ICU Backend: Services, State, and Alerts}
%==============================================================================

\section{FastAPI application}
The main app mounts ICU routers under \texttt{/api/icu}, exposes training and dataset utilities under \texttt{/api/...}, and enables CORS for localhost and private LAN patterns so phones and tablets can load the SPA while calling the API.

\section{Representative ICU endpoints}
\begin{table}[H]
\centering
\caption{Selected ICU REST capabilities (illustrative; see source for full list).}
\begin{tabular}{@{}p{6.8cm}p{7.8cm}@{}}
\toprule
\textbf{Route (prefix \texttt{/api/icu})} & \textbf{Purpose} \\
\midrule
\texttt{POST /patient} & Create or update patient metadata for admission \\
\texttt{POST /bed/\{id\}/assign-patient} & Bind patient to bed \\
\texttt{POST /bed/\{id\}/vitals} & Ingest HR, temperature, SpO$_2$, BP, RR, etc. \\
\texttt{POST /bed/\{id\}/classification} & Ingest neuromorphic beat decisions and MAC readout \\
\texttt{POST /bed/\{id\}/ecg} & Ingest ECG waveform segments for plotting \\
\texttt{GET /summary} & Ward summary for dashboards \\
\bottomrule
\end{tabular}
\end{table}

\section{WebSocket summaries}
The \texttt{/ws/icu} socket emits periodic JSON summaries so the UI stays synchronized without polling every REST resource. Event-driven optimizations remain possible while preserving a simple polling loop internally.

\section{Alert engine}
Threshold rules catch tachycardia, bradycardia, fever, and related vitals excursions; trend rules consider sustained patterns (e.g., prolonged elevation or accumulating ectopy). Alerts surface in the UI with severities suitable for nursing and physician workflows.

\section{Mortality risk model}
A joblib-packaged model under \texttt{ICU/models/} exposes feature-ordered inference via \texttt{predict\_mortality\_risk}, integrating structured inputs into a probability used by the clinical UI for early warning-style displays. This augments beat-level neuromorphic outputs with aggregate risk context.

%==============================================================================
\chapter{Frontend: Ward Dashboard, Bed Detail, and Mobile Use}
%==============================================================================

\section{Responsive ICU user experience}
The React application provides:
\begin{itemize}[itemsep=0.2em]
  \item \textbf{Multi-bed ward overview} with connection status, headline vitals, and highest-severity alert per bed.
  \item \textbf{Per-bed detail} with live ECG strip, classification stream, vitals trends, orders, alert response, and doctor notes.
  \item \textbf{Admission flows} capturing demographics and clinical context to support continuity of care.
\end{itemize}

\section{Mobile access for physicians}
The startup script binds Vite and Uvicorn to all interfaces and prints \textbf{LAN URLs} (e.g., \texttt{http://192.168.x.x:5173}) so clinicians on the same Wi-Fi open the ward view from phones or tablets without VPN complexity in lab deployments. Environment variables \texttt{VITE\_API\_URL} and \texttt{VITE\_WS\_URL} point to the machine's LAN IP, ensuring browser calls hit the correct backend from mobile clients.

\subsection{Advantages of mobile ward access}
\begin{itemize}[itemsep=0.2em]
  \item \textbf{Rapid situational awareness} when moving between rooms or during procedures outside fixed workstation desks.
  \item \textbf{Shared team visibility} on rounds, with consistent data to attending and residents.
  \item \textbf{Alert triage} from pocket devices, reducing time-to-acknowledgment when paired with appropriate hospital policy and authentication.
\end{itemize}

\section{User interface figures}
\begin{figure}[H]
  \centering
  \fig[width=0.95\textwidth]{Admin\ Auth.png}
  \caption{Administrative authentication screen: gated access to ward operations and patient-identifying views.}
  \label{fig:auth}
\end{figure}

\begin{figure}[H]
  \centering
  \fig[width=0.95\textwidth]{NEW\ ADMISSION.png}
  \caption{New admission workflow: structured capture of patient identifiers and clinical metadata tied to bed assignment.}
  \label{fig:admission}
\end{figure}

\begin{figure}[H]
  \centering
  \fig[width=0.95\textwidth]{Full\ ICU\ bed\ dashboard.png}
  \caption{Full ICU bed dashboard (desktop): live ECG, vitals, neuromorphic/ML panels, and operational controls in one view.}
  \label{fig:dashboard-desktop}
\end{figure}

\begin{figure}[H]
  \centering
  \fig[width=0.85\textwidth]{Full\ ICU\ bed\ dashboard\ mobile.jpeg}
  \caption{Bed dashboard on a mobile form factor: responsive layout preserving critical waveforms and alerts for physicians on the move.}
  \label{fig:dashboard-mobile}
\end{figure}

\begin{figure}[H]
  \centering
  \fig[width=0.95\textwidth]{Bed\ 2\ Dashboard.png}
  \caption{Secondary bed view illustrating scalable multi-bed monitoring and consistent visual language across beds.}
  \label{fig:bed2}
\end{figure}

\begin{figure}[H]
  \centering
  \fig[width=0.85\textwidth]{Bed\ 2\ dashboard\ moble\ mode.jpeg}
  \caption{Mobile-oriented layout for an additional bed, reinforcing LAN-based access patterns used in ward walkthroughs.}
  \label{fig:bed2mobile}
\end{figure}

\begin{figure}[H]
  \centering
  \fig[width=0.95\textwidth]{ML\ prediction.png}
  \caption{Machine-learning prediction panel: combines global and patient-contextual outputs for decision support alongside raw signals.}
  \label{fig:ml}
\end{figure}

\begin{figure}[H]
  \centering
  \fig[width=0.95\textwidth]{Doctors\ notes.png}
  \caption{Physician notes interface: free-text documentation linked to the active encounter and synchronized with ward state.}
  \label{fig:notes}
\end{figure}

\FloatBarrier

%==============================================================================
\chapter{Deployment, Networking, and Bedside Operations}
%==============================================================================

\section{Unified startup (\texttt{start.sh})}
The repository provides a single shell entry point that prepares a consistent developer and ward-station environment. On launch, the script terminates stray processes from prior sessions (Uvicorn, \texttt{serial\_bridge.py}, Vite), clears Python \texttt{\_\_pycache\_\_} caches and frontend build caches, then starts three cooperating services: the FastAPI backend on \texttt{0.0.0.0:8000}, the optional serial bridge for USB-connected ESP boards, and the Vite development server bound to all interfaces so tablets and phones on the LAN can attach.

\section{LAN discovery and mobile pairing}
On macOS the script queries \texttt{ipconfig} for \texttt{en0}/\texttt{en1}; on Linux it uses \texttt{hostname -I}. When a private IPv4 address is found, it exports \texttt{VITE\_API\_URL=http://<LAN>:8000/api} and \texttt{VITE\_WS\_URL=ws://<LAN>:8000} before starting the frontend. This is the critical link between \textbf{authentic bedside HTTP posts} from controllers and \textbf{physician phones}: the browser on a mobile device must address the backend by the same routable IP, not \texttt{localhost}, which would refer to the phone itself.

\section{CORS and same-network policy}
The backend enables CORS for common localhost ports and, via regular expression, RFC1918 private ranges. Together with explicit binding to \texttt{0.0.0.0}, the stack supports ward trials where clinicians type a single printed URL on their handsets without reconfiguring each build artifact.

\section{Serial bridge}
For laboratory benches, \texttt{backend/serial\_bridge.py} can forward UART telemetry from development kits into the same logical monitoring path as Wi-Fi POSTs, preserving a single ingestion contract when tethered and wireless controllers are used together during bring-up.

%==============================================================================
\chapter{Clinical Workflows Supported by the Stack}
%==============================================================================

\section{Admission and identity}
Creating a patient via \texttt{POST /api/icu/patient} establishes identifiers and structured metadata (age, gender, diagnosis, contact, medical record number, admission date). Assigning the patient to a bed couples the logical person to a physical location for telemetry routing. The UI flows in Figures~\ref{fig:auth}--\ref{fig:admission} mirror this sequence: authentication first, then structured intake so every subsequent sample is tagged to a real encounter rather than an anonymous stream.

\section{Rounds and continuous review}
The ward overview aggregates each bed's connectivity, headline vitals, and worst active alert. During rounds, a team can start from the multi-bed tile view, drill into a detail page with live ECG and classification history, and cross-check peripheral pulse rate against ECG-derived rate. Doctor notes (Figure~\ref{fig:notes}) capture contemporaneous reasoning, orders panels support actionable items, and discharge portals close the loop when the patient leaves the unit---all within one coherent React application.

\section{Mobile-first advantages in practice}
Clinicians are rarely stationary. The mobile layouts (Figures~\ref{fig:dashboard-mobile} and~\ref{fig:bed2mobile}) demonstrate that waveform-centric monitoring need not be confined to central nursing stations. Benefits include faster recognition of arrhythmia burden changes when a physician is already en route to the room, parallel awareness during sterile procedures where a desktop is unavailable, and teaching moments where trainees mirror the attending's live screen on personal devices under supervision.

%==============================================================================
\chapter{Advantages of the Integrated Neuromorphic ICU Platform}
%==============================================================================

\section{Edge intelligence and latency}
Ventricular ectopy and related beat-level phenomena are time-sensitive. Performing MAC inference in analog hardware after spike encoding avoids shipping every raw sample to a data center for GPU batching. The neuromorphic path keeps classification aligned with the physical timing of depolarization, while the backend still receives downsampled segments for charting and longitudinal analytics.

\section{Energy and cost profile}
Student and research labs can replicate the resistive MAC with inexpensive SPI potentiometers and commodity op-amps. Compared with always-on cloud inference, edge analog MAC reduces recurring compute charges and respects air-gapped hospital segments that cannot export PHI externally.

\section{Openness and extensibility}
Every layer---schematics, firmware, Python training, FastAPI routes, React components---is available for inspection. Hospitals can adapt alert thresholds, add sensors, or swap classifiers without vendor lock-in, subject to local regulatory processes.

\section{Multi-sensor redundancy}
Pulse and temperature provide orthogonal checks. Fever plus tachycardia may trigger different escalation than either alone; ECG ectopy burden adds arrhythmic context on top of rate scalars.

\section{Personalization without data mixing}
Digital twins fine-tune from one patient's longitudinal CSV and event history, improving fit to morphology while the training pipeline enforces strict \texttt{patient\_id} filters.

%==============================================================================
\chapter{Frontend Architecture and Key Components}
%==============================================================================

\section{Application structure}
The SPA uses stage-based pages (authentication, ward, detail) and shared components. Representative modules include \texttt{MultiBedDashboard} for ward tiles, \texttt{RealTimeEcgMonitor} for scrolling traces, \texttt{EarlyPredictionPanel} and ML surfaces (Figure~\ref{fig:ml}), \texttt{AlertResponsePanel} for acknowledging escalations, \texttt{OrdersEditPanel} for interventions, \texttt{DoctorNotesPanel} for narrative documentation, and admission/discharge portals for encounter boundaries. Styling splits into co-located CSS modules to preserve readability as features accrue.

\section{Real-time transport}
Initial summaries arrive via REST; \texttt{/ws/icu} pushes periodic JSON snapshots so tiles refresh without aggressive polling. Client code resolves API bases from \texttt{import.meta.env} so production builds can target institutional hostnames.

%==============================================================================
\chapter{Backend Internals: State, Logging, and Risk}
%==============================================================================

\section{\texttt{HardwareState} and per-bed aggregation}
Each bed maintains recent vitals samples, rolling ECG buffers, classification timelines, and alert sets. Controllers may post at different cadences; the server merges by timestamps to present a coherent timeline to the UI.

\section{Patient-associated CSV logs}
For supported configurations, vitals append to per-patient files under \texttt{data/icu\_logs/}. Rows include ISO timestamps, identifiers, heart rate, temperature, SpO$_2$, blood pressure components, respiratory rate, and fall-detection flags when provided. These files ground digital twin training and offline audits; they are ordinary CSV text, readable in spreadsheets or pandas without proprietary tooling.

\section{Mortality risk scoring}
The joblib artifact in \texttt{ICU/models/} stores a trained classifier and ordered feature names. \texttt{predict\_mortality\_risk} builds a dense feature row with zeros for missing entries, then returns the positive-class probability from \texttt{predict\_proba}. The UI can surface this alongside beat labels to combine \textbf{acute rhythm information} with \textbf{aggregate risk} in one glance.

%==============================================================================
\chapter{Validation, Calibration, and Bench Procedures}
%==============================================================================

\section{Comparator thresholding}
Spike density depends on threshold selection relative to noise floor. Bench validation correlates serial-printed ECG amplitude with oscilloscope captures to ensure the LM393 toggles once per QRS complex under nominal gain.

\section{Memristor map verification}
After training, exported resistances should be spot-checked with a multimeter or bridge measurement on the wiper nodes. Firmware maps floating-point ohms to digital pot codes; saturation at extremes indicates the need to rescale weights or adjust reference voltages.

\section{Wi-Fi stress}
High packet loss manifests as gaps in plotted waveforms. For ICU prototypes, place access points with clear line-of-sight and separate SSIDs for patient traffic versus guest networks.

%==============================================================================
\chapter{Comparison with Conventional Monitor Architectures}
%==============================================================================

\begin{table}[H]
\centering
\caption{Qualitative comparison: integrated neuromorphic research platform vs typical proprietary bedside monitors.}
\begin{tabular}{@{}p{4.2cm}p{5.4cm}p{5.4cm}@{}}
\toprule
\textbf{Aspect} & \textbf{Commercial ICU monitor} & \textbf{This platform} \\
\midrule
Algorithm transparency & Black-box vendor DSP & Open Python training + inspectable weights \\
Personalization & Limited or cloud-only & Per-patient twins with optional HW map \\
Edge inference style & ASIC DSP & Memristor-emulating analog MAC \\
Extensibility & Vendor timelines & Fork firmware and backend freely \\
Mobile access & Often separate product & Same React app via LAN URL \\
\bottomrule
\end{tabular}
\end{table}

%==============================================================================
\chapter{Digital Twins and the Super Model}
%==============================================================================

\section{Super model}
The population-trained linear classifier (and exported memristor map) acts as the \textbf{super model}: a stable default when patient-specific adaptation is not yet available.

\section{Per-patient digital twins}
When enough patient-tagged beats and vitals accumulate, a twin training job can initialize from super-model weights and refine using only that patient's history. Isolation guarantees prevent cross-patient leakage; optional per-patient JSON maps can be deployed to hardware for personalized neuromorphic inference at the bedside.

%==============================================================================
\chapter{Security, Privacy, and Operational Considerations}
%==============================================================================

\section{Network posture}
LAN deployment with explicit CORS rules trades Internet-wide exposure for controlled clinical engineering environments. Production deployments should add TLS termination, VPN or Zero Trust access, and institutional identity providers---outside the student-lab scope but architecturally compatible.

\section{Patient data isolation}
Every persisted structure is keyed by \texttt{patient\_id}; training filters enforce separation. CSV logs use per-patient filenames for human-auditable organization.

%==============================================================================
\chapter{Performance and Engineering Metrics}
%==============================================================================

\section{Latency}
Edge inference paths target sub-millisecond MAC evaluation after window completion; interrupt overhead remains bounded by spike rate and window length. Backend ingestion is lightweight JSON parsing and state updates suitable for tens of beds on modest hardware in prototype configurations.

\section{Power}
Bench measurements on ESP32-class boards typically fall in the hundreds of mA active depending on Wi-Fi usage; duty-cycled designs can reduce average power when radios are quieted between posts.

%==============================================================================
\chapter{Future Work}
%==============================================================================

\begin{itemize}[itemsep=0.25em]
  \item Expand sensor suite with certified SpO$_2$ and NIBP modules; integrate capnography for respiratory failure surveillance.
  \item Harden WebSocket pushes to event-driven alert notifications; add authenticated roles (nurse, physician, admin).
  \item Formal clinical validation study with ethics approval, independent of this engineering report.
  \item Explore memristor devices beyond potentiometers where fabrication partnerships allow.
\end{itemize}

%==============================================================================
\chapter{Expanded Discussion: Nonidealities and Mitigations}
%==============================================================================

\section{Electrode impedance and motion}
Skin--electrode impedance varies with preparation and patient diaphoresis. High impedance increases thermal noise and motion artifact, distorting both ADC traces and comparator crossings. Standard mitigation includes abrasive prep where clinically appropriate, fresh gelled electrodes, and mechanical strain relief on lead wires. Software-side, the training pipeline's bandpass and robust normalization reduce sensitivity to slow baseline shifts that survive analog high-passing.

\section{Comparator chatter and hysteresis}
Comparators can chatter when inputs dwell near the threshold. Small hysteresis networks or RC filtering on the input pin may be required in noisy environments. Firmware refractory windows suppress double-counting within a single QRS when edges remain clean.

\section{Quantization and ADC limits}
ESP32 ADCs are not laboratory digitizers; absolute voltage accuracy matters less than repeatability for relative MAC comparisons when bias is calibrated jointly. For research-grade archiving, consider external sigma-delta ADCs on SPI while retaining the neuromorphic path for inference.

\section{Wi-Fi contention}
802.11 shared medium behavior introduces jitter in HTTP arrival times. Waveform plots should tolerate small gaps; alert rules should rely on sustained excursions rather than single missing packets unless connectivity loss itself is an alert category.

\section{Thermal drift in analog components}
Resistor networks and op-amp offsets drift with temperature. Bench validation should include warm-up periods; production deployments may schedule periodic bias recalibration against a known test pattern.

%==============================================================================
\chapter{Team, Provenance, and Event Context}
%==============================================================================

This system was engineered by \textbf{Team BitWizards} as part of the \textbf{Visualize'26} hackathon, integrating open scientific corpora for classifier training with original hardware and software integration work described herein. The repository README cites PhysioNet/MIT-BIH and community tooling; this report extends that summary into a single narrative suitable for coursework, portfolio review, or sponsor briefings.

%==============================================================================
\chapter{Conclusion}
%==============================================================================

This project delivers a \textbf{working, open, end-to-end path} from custom analog neuromorphic ECG acquisition through multi-sensor ICU fusion, AI-assisted visualization, mobile-ready ward software, and optional personalized twins. The hardware is real and testable on the bench; the software stack ingests authenticated bedside streams; the photographs document both the physical system and the operational UI. Together, these elements form a coherent platform for research and teaching in neuromorphic edge intelligence applied to critical care.

%==============================================================================
\appendix
%==============================================================================

\chapter{Mathematical Formulations: Preprocessing and Classification}
%==============================================================================

\section{Notation}
Let $f_s$ denote sampling rate in Hz, $N$ the number of samples in a window, and $\mathbf{x}\in\mathbb{R}^N$ a single beat-aligned excerpt after preprocessing. Spike vectors $\mathbf{s}\in\{0,1\}^N$ feed the linear model equivalently to $\mathbf{x}$ when training and hardware paths are kept consistent.

\section{Baseline wander removal}
High-pass filtering removes DC and low-frequency drift. The normalized cutoff is $f_{\text{norm}} = f_c/(f_s/2)$ with Nyquist $f_s/2$. For a 0.5~Hz high-pass at $f_s=360$~Hz, $f_{\text{norm}}=1/f_s$ in normalized frequency units used by SciPy routines. Butterworth order is typically four; zero-phase (\texttt{filtfilt}) application avoids phase distortion that would misalign R-peaks.

\section{Bandpass filtering}
Butterworth bandpass with $f_L = 0.5$~Hz and $f_H = 40$~Hz retains QRS energy while suppressing baseline and high-frequency noise:
\begin{equation}
  f_{L,\text{norm}} = \frac{f_L}{f_s/2}, \qquad f_{H,\text{norm}} = \frac{f_H}{f_s/2}.
\end{equation}

\section{Normalization}
Min--max scaling maps samples to $[0,1]$:
\begin{equation}
  x' = \frac{x - x_{\min}}{x_{\max} - x_{\min}}, \qquad \tilde{x} = x'(b-a) + a.
\end{equation}
If $x_{\max}=x_{\min}$, the segment is constant and mapped to the lower bound to avoid division by zero. Z-score,
\begin{equation}
  \tilde{x} = \frac{x-\mu}{\sigma},
\end{equation}
and robust median/IQR scaling are supported when outliers dominate.

\section{Beat segmentation}
For R-peak index $n_R$, windows of length 216 samples (600~ms at 360~Hz) use offsets (e.g., 72 samples pre-peak, 144 post) yielding $\mathbf{x} \in \mathbb{R}^{216}$. The asymmetric split emphasizes post-QRS ST/T segments relevant to ventricular ectopy discrimination.

\section{Spike encoding}
\subsection{Threshold encoding}
\begin{align}
  s_i &= 1 \quad \text{if } x_i > \theta \quad (\text{``above'' mode}), \\
  s_i &= 1 \quad \text{if } x_i < \theta \quad (\text{``below'' mode}); \quad \text{else } s_i = 0.
\end{align}
Refractory suppression discards spikes closer than $n_{\text{ref}}$ samples.

\subsection{Rate encoding}
Spike probability per sample can follow
\begin{equation}
  P(\text{spike at } i) = \tilde{x}_i \, r_{\max} \, \Delta t, \qquad \Delta t = \frac{1}{f_s},
\end{equation}
with $\tilde{x}$ normalized to $[0,1]$ and $r_{\max}$ the maximum firing rate (e.g., 100~Hz).

\section{Linear classifier}
\begin{equation}
  z = \mathbf{w}^\top \mathbf{x} + b, \qquad \hat{y} = \mathbb{I}[z > 0].
\end{equation}
For logistic regression,
\begin{equation}
  P(y{=}1\mid \mathbf{x}) = \sigma(z) = \frac{1}{1+e^{-z}}.
\end{equation}
Training minimizes regularized cross-entropy; hardware compares MAC voltage to a bias threshold calibrated from $b$.

\section{Weight-to-resistance mapping}
Linear mapping:
\begin{equation}
  R_i = R_{\min} + \frac{w_i - w_{\min}}{w_{\max} - w_{\min}}(R_{\max} - R_{\min}).
\end{equation}
Symmetric mapping for signed weights uses $\hat{w}_i = w_i/\max_j|w_j|$ and
\begin{equation}
  R_i = R_{\text{center}} - \hat{w}_i \, \Delta R,
\end{equation}
clipped to $[R_{\min},R_{\max}]$. Conductance mapping sets $G_i = G_{\min}+\hat{w}_i(G_{\max}-G_{\min})$ with $R_i=1/G_i$.

\section{Evaluation metrics}
\begin{align}
  \text{Accuracy} &= \frac{\text{TP} + \text{TN}}{\text{TP} + \text{TN} + \text{FP} + \text{FN}}, \\
  \text{Sensitivity} &= \frac{\text{TP}}{\text{TP} + \text{FN}}, \qquad
  \text{Specificity} = \frac{\text{TN}}{\text{TN} + \text{FP}}, \\
  \text{Precision} &= \frac{\text{TP}}{\text{TP} + \text{FP}}, \\
  F_1 &= \frac{2\,\text{Precision}\,\text{Recall}}{\text{Precision} + \text{Recall}}.
\end{align}

\section{Analog MAC relation}
For reference $V_{\text{ref}}$ and selected resistors $R_i$ with feedback $R_f$,
\begin{equation}
  V_{\text{MAC}} = -R_f \sum_i \frac{V_{\text{ref}}}{R_i} \, s_i,
\end{equation}
with $s_i \in \{0,1\}$ the switch state driven by spike encoding.

\section{Divider-based resistance inference}
For series $R$ and $R_{\text{ref}}$ driven by $V_{\text{supply}}$, midpoint voltage $V_m$ yields
\begin{equation}
  R = R_{\text{ref}}\left(\frac{V_{\text{supply}}}{V_m} - 1\right).
\end{equation}

\chapter{Step-by-Step Training and Export Workflow}
%==============================================================================

\begin{enumerate}[itemsep=0.35em]
  \item Acquire MIT-BIH records into \texttt{data/mitbih} using the dataset downloader.
  \item Parse waveforms and annotations; retain beats labeled \texttt{N} and \texttt{V}.
  \item Apply high-pass, bandpass, and normalization per \texttt{global\_config.yaml}.
  \item Segment fixed windows around annotated R-peaks.
  \item Encode windows to spike trains (\texttt{threshold} or \texttt{rate}).
  \item Train logistic regression or SGD linear classifiers; extract $\mathbf{w}$ and $b$.
  \item Map weights to resistance targets within hardware bounds.
  \item Emit \texttt{memristor\_map.json} including thresholds, refractory period, and metadata (accuracy, date, record list).
  \item Upload JSON to ESP32 SPIFFS; flash firmware; verify spikes and MAC on serial.
  \item Enable ICU HTTP client firmware variant to stream classifications and ECG segments to FastAPI.
\end{enumerate}

\chapter{Firmware Pin Map (ESP32-WROOM-32D Example)}
%==============================================================================

\begin{table}[H]
\centering
\caption{Indicative GPIO assignments from project firmware notes (verify against \texttt{config.h}).}
\begin{tabular}{@{}ll@{}}
\toprule
\textbf{Signal} & \textbf{GPIO} \\
\midrule
ECG ADC input & 34 \\
Comparator spike input & 4 (interrupt) \\
MAC output ADC & 35 \\
MCP4131 chip-select lines & 5, 15 \\
CD4066 control lines & 25, 26 \\
\bottomrule
\end{tabular}
\end{table}

\chapter{CSV Vitals Schema (Bed~1 Logging)}
%==============================================================================

Logged CSV columns include: \texttt{timestamp}, \texttt{bed\_id}, \texttt{patient\_id}, \texttt{heart\_rate}, \texttt{temperature}, \texttt{spo2}, \texttt{bp\_systolic}, \texttt{bp\_diastolic}, \texttt{resp\_rate}, \texttt{fall\_detected}. Filenames follow \texttt{bed1\_<patient\_id>\_<sanitized\_name>.csv}, supporting traceable per-patient archives for twin training and quality review.

\chapter{Configuration Reference (Excerpt)}
%==============================================================================

Key Python paths appear in \texttt{config/global\_config.yaml}: dataset directories, allowed labels \texttt{N}/\texttt{V}, filter bands, normalization mode, spike parameters, classifier regularization, train/validation splits, and memristor bounds. Firmware mirrors window sizes and GPIO mappings in \texttt{esp32\_firmware/src/config.h}. Memristor-specific scaling appears in \texttt{config/memristor\_specs.yaml}.

\chapter{Bill of Materials and Bench Roles (Indicative)}
%==============================================================================

\begin{table}[H]
\centering
\caption{Representative components for the neuromorphic ECG path.}
\begin{tabular}{@{}ll@{}}
\toprule
\textbf{Component} & \textbf{Role} \\
\midrule
ESP32-S3 / ESP32-WROOM & Control, ADC, Wi-Fi, interrupts \\
AD620 & Instrumentation amplification \\
LM393 & Comparator spike generation \\
CD4066 & Analog switch matrix \\
MCP4131 / AD5204 & Programmable resistances (memristor emulation) \\
OPA227 / LM358 & Filtering, MAC summing, buffering \\
N-channel MOSFET & Optional pulse shaping for programming \\
\bottomrule
\end{tabular}
\end{table}

\begin{table}[H]
\centering
\caption{Auxiliary vitals path (typical).}
\begin{tabular}{@{}ll@{}}
\toprule
\textbf{Component} & \textbf{Role} \\
\midrule
ESP8266 module & ADC for pulse/temperature, Wi-Fi client \\
Optical pulse sensor & Peripheral HR waveform \\
LM35 & Body temperature voltage \\
\bottomrule
\end{tabular}
\end{table}

\chapter{Operational Checklists (Bench and Ward)}
%==============================================================================

\section{Hardware bring-up}
\begin{itemize}[itemsep=0.2em]
  \item Verify supply rails within absolute maximum ratings for all ICs.
  \item Confirm AD620 reference and sense lines match electrode cable pinout.
  \item Sweep comparator threshold while injecting synthetic sine to confirm single pulses per cycle.
  \item Program pots to mid-scale before attaching high gain; avoid saturating summing nodes.
  \item Confirm SPI chip-select uniqueness across MCP4131/AD5204 devices.
\end{itemize}

\section{Software and networking}
\begin{itemize}[itemsep=0.2em]
  \item Run backend with \texttt{0.0.0.0} binding; open firewall ports only on trusted VLANs.
  \item Print LAN URL from \texttt{start.sh}; verify mobile device reaches \texttt{/api} root JSON.
  \item Confirm WebSocket summary messages arrive in browser devtools network tab.
  \item Rotate any shared credentials used in ward pilots; prefer institutional SSO in production.
\end{itemize}

\section{Data stewardship}
\begin{itemize}[itemsep=0.2em]
  \item Restrict CSV directories to clinical roles; encrypt disks on laptops holding logs.
  \item Document firmware version and memristor map hash alongside exported vitals for reproducibility.
\end{itemize}

\chapter{Glossary}
%==============================================================================

\begin{description}[style=nextline]
  \item[ECG] Electrocardiogram: electrical potential time series reflecting cardiac depolarization/repolarization.
  \item[MAC] Multiply--accumulate: here implemented as summed currents through weighted resistors.
  \item[Super model] Population-trained classifier weights used as the default and twin initialization.
  \item[Digital twin] Patient-specific model fine-tuned only on that patient's data.
  \item[SAM fusion] Project term for synchronized aggregation of multi-sensor vitals at the bed.
  \item[LAN URL] Host address reachable from phones on the same Wi-Fi as the backend (not loopback).
  \item[CORS] Cross-Origin Resource Sharing policy governing browser access to APIs on different origins.
\end{description}

\begin{thebibliography}{99}
\bibitem{mitbih} Moody, G. B., Mark, R. G. MIT-BIH Arrhythmia Database. PhysioNet. (Goldberger et al., contributors).
\bibitem{physionet} Goldberger, A. L., et al. PhysioBank, PhysioToolkit, and PhysioNet: Components of a New Research Resource for Complex Physiologic Signals. Circulation 101(23):e215--e220 (2000).
\bibitem{esp-idf} Espressif Systems. ESP32 Technical Reference Manual and Arduino-ESP32 Core Documentation.
\bibitem{ad620} Analog Devices. AD620: Low Cost, Low Power Instrumentation Amplifier. Datasheet.
\bibitem{lm393} Texas Instruments. LM393/LM293/LM193 Dual Differential Comparators. Datasheet.
\bibitem{cd4066} Texas Instruments. CD4066B CMOS Quad Bilateral Switch. Datasheet.
\bibitem{mcp4131} Microchip. MCP4131 Single 128-Taps Digital Potentiometer with SPI. Datasheet.
\bibitem{opa227} Texas Instruments. OPA227/OPA228 Ultra-Low Noise Precision Operational Amplifiers. Datasheet.
\bibitem{fastapi} FastAPI Documentation. \url{https://fastapi.tiangolo.com/}
\bibitem{react} Meta Open Source. React---A JavaScript library for building user interfaces. \url{https://react.dev/}
\bibitem{sklearn} Pedregosa, F., et al. Scikit-learn: Machine Learning in Python. JMLR 12:2825--2830 (2011).
\bibitem{wfdb} Silva, I., et al. The PhysioNet/Computing in Cardiology Challenge 2013: Noninvasive Fetal ECG. Computers in Cardiology 40:137--140 (2013). (wfdb tooling ecosystem.)
\bibitem{ieee1073} ISO/IEEE 11073 Family (overview of medical device communication; contextual for future interoperability work).
\bibitem{hl7fhir} HL7 FHIR Specification (R4+). \url{https://www.hl7.org/fhir/} (reference for future hospital integration).
\bibitem{neuromorphic} Indiveri, G., et al. Integration of nanoscale memristor synapses in neuromorphic computing architectures. Nanotechnology 24(38):384010 (2013).
\bibitem{icu_ml} Johnson, A. E. W., et al. MIMIC-III, a freely accessible critical care database. Scientific Data 3:160035 (2016). (context for ICU ML, separate from this hardware build.)
\end{thebibliography}

\end{document}
